\section*{Schedules}
\label{sec:schedules}

\scheduleSection{Uniform convergence}{4}{
The general principle of uniform convergence. A uniform limit of continuous functions is continuous. Limits of differentiable functions are differentiable provided the derivatives (not assumed continuous) converge uniformly. Series of functions, the Weierstrass M-test. Local uniform convergence of power series.}

\scheduleSection{Uniform continuity and integration}{3}{Uniform continuity. Continuous functions on closed bounded intervals are uniformly continuous. Riemann integration, including review of basic facts (from Analysis I). A uniform limit of integrable functions is integrable, and the integral of the limit is the limit of the integrals.}

\scheduleSection{Metric spaces}{6}{Definition of a metric space, examples including the Euclidean metric on \(\RR^n\). Limits, open and closed sets. Sequential compactness, Bolzano--Weierstrass theorem. Continuity, every continuous function on sequentially compact set is bounded and is uniformly continuous. Definition of a norm; examples, including matrix norms; all norms on finite-dimensional vector spaces are (Lipschitz) equivalent. Path-connectedness.}

\scheduleSection{Differentiation from \(\RR^m\) to \(\RR^n\)}{6}{
Definition of derivative as a linear map; elementary properties; the chain rule. Partial derivatives; continuous partial derivatives imply differentiability. Higher-order derivatives; symmetry of mixed partialderivatives (assumed con tinuous). Hessian, second order Taylor expansion for real-valued \(\contClass{2}\) functions, classification of critical points. The mean-value inequality. A function having zero derivative on a path-connected open subset is constant.}

\scheduleSection{The contraction mapping theorem}{3}{Definition of completeness, examples including incomplete spaces. Completeness of \(\RR^n\) and \(\contClass{0}\pqty{\ccintv{a}{b}}\) with the uniform norm; a subset of a complete space is complete if and only if it is closed. The contraction mapping theorem.}

\scheduleSection{Applications of the contraction mapping theorem}{2}{The inverse function theorem; the implicit function theorem. Picard's solution of systems of ordinary differential equations.}